% ===================================================================
%  QUADRO N.º 9 — A Montanha. — A Estância termal. — A Floresta.
%
%  Ceci n'est PAS une transcription : c'est une traduction, faite en
%  2026, en portugais d'aujourd'hui. Voir texto/pt/10-quadro-01.tex
%  pour la consigne, qui vaut pour les seize fichiers.
% ===================================================================

%%K t09-apar-1 apar
\VUsaut{4.0mm}
\VUtitre{15.0pt}{60}{QUADRO N.º 9}
\VUsaut{3.0mm}

%%K t09-tit-1 sub
\VUcentre{12.0pt}{0}{\textit{Primeira cena.}}
\VUsaut{3.0mm}

%%K t09-tit-2 sub
\VUpk{11.4pt}{40}{A Montanha. --- A Estância termal.}
\VUsaut{3.0mm}

%%K t09-c1-01-1 p
1. --- Estou há oito dias em Pratz. Não consigo exprimir toda a admiração
que senti quando me encontrei diante da maravilhosa \VUgras{cordilheira}
\textsuperscript{(1)} dos Pirenéus, cuja \VUgras{crista} \textsuperscript{(2)} se
recorta sobre o fundo azul do céu como um festão gigantesco.

%%K t09-c1-02-1 p
2. --- Não imaginava que os \VUgras{cumes} \textsuperscript{(3)} pirenaicos
alcançassem uma \VUgras{altura} tão grande, e fiquei pasmado diante dos
magníficos \VUgras{glaciares} \textsuperscript{(5)} e dos \VUgras{picos}
\textsuperscript{(4)} cobertos de neve pelos quais rolam \VUgras{avalanches} tão
terríveis.

%%K t09-tit-3 sub
\VUsaut{3.0mm}
\VUpk{10.2pt}{40}{Paisagem de montanha.}
\VUsaut{3.0mm}

%%K t09-c1-03-1 p
3. --- No dia seguinte mesmo à minha chegada fui ao velho \VUgras{vulcão}
\textsuperscript{(6)} apagado que há perto de Pratz. Nas bordas áridas e nuas
da \VUgras{cratera} veem-se ainda muito bem os vestígios da passagem da
\VUgras{lava} que corria, há vários séculos, pela \VUgras{encosta da
montanha} \textsuperscript{(7)}. Consegui encontrar, numa das \VUgras{vertentes},
alguns fragmentos dessa lava.

%%K t09-c1-04-1 p
4. --- Pratz está na fronteira, e a \VUgras{fortaleza} \textsuperscript{(8)} que
defende a \VUgras{garganta} \textsuperscript{(27)} que separa a França da Espanha
lembra em tudo aqueles velhos castelos das margens do Reno que parecem
espreitar uma \VUgras{presa} do alto do seu rochedo.

%%K t09-c1-05-1 p
5. --- Uma \VUgras{águia} \textsuperscript{(9)} enorme, que tem o seu ninho não
longe do \VUgras{forte} \textsuperscript{(8)}, atrai muitas vezes a minha atenção.
Esta \VUgras{ave de rapina}, de \VUgras{bico} poderoso, leva com frequência
às suas crias um cordeiro ou um cabrito preso nas suas \VUgras{garras}. As
suas \VUgras{asas} são tão grandes que pode planar longo tempo com a sua
pesada carga.

%%K t09-tit-4 sub
\VUsaut{3.0mm}
\VUpk{10.2pt}{40}{O Excursionista.}
\VUsaut{3.0mm}

%%K t09-c1-06-1 p
6. --- O passeio mais agradável daqui é a excursão à \VUgras{cascata}
\textsuperscript{(10)} de Cau. A \VUgras{queda de água} cai em remoinhos e
depois perde-se num bonito \VUgras{ribeiro} no \VUgras{pinhal de abetos}
\textsuperscript{(11)} situado a oeste da cidade. Subimos por esse bosque até ao
\VUgras{observatório meteorológico} \textsuperscript{(12)} e, do alto do
\VUgras{miradouro}, abarcamos o \VUgras{panorama} inteiro da região. A
\VUgras{paisagem} é soberba e a \VUgras{natureza} parece ter prodigalizado
a este país todos os tesouros de que dispõe.

%%K t09-c1-07-1 p
7. --- Para descer tomámos o \VUgras{funicular} \textsuperscript{(13)} e parámos
numa pequena \VUgras{estação} junto às \VUgras{ruínas} \textsuperscript{(14)} de um
velho \VUgras{castelo feudal} habitado outrora pelos senhores de Pratz.

%%K t09-c1-08-1 p
8. --- Pobre castelo! Que pensará ao ver a seus pés a coquete
\VUgras{cidade balnear} \textsuperscript{(15)}, que hoje é uma \VUgras{estância
termal} da moda?

%%K t09-c1-08-2 p
O \VUgras{clima} é tão agradável e a \VUgras{água mineral} tão benéfica que
a \VUgras{cura} que se segue no \VUgras{sanatório} \textsuperscript{(17)} é um
verdadeiro prazer.

%%K t09-tit-5 sub
\VUsaut{3.0mm}
\VUpk{10.2pt}{40}{Uma agradável cidade termal.}
\VUsaut{3.0mm}

%%K t09-c1-09-1 p
9. --- De resto, fez-se tudo para que a estada seja agradável.
Construiu-se um grande \VUgras{viaduto} \textsuperscript{(16)} para se poder
lançar uma via férrea; o \VUgras{parque} \textsuperscript{(18)} das
\VUgras{termas}, no qual se dão \VUgras{concertos}, está arranjado de um
modo encantador. Construíram-se vários graciosos \VUgras{chalés}
\textsuperscript{(19)} à volta do \VUgras{casino} \textsuperscript{(21)}, e uma
\VUgras{calçada} \textsuperscript{(22)} muito bem tratada torna facílimas as
comunicações.

%%K t09-c1-10-1 p
10. --- Um simples \VUgras{talude} \textsuperscript{(23)} separa a \VUgras{estrada}
\textsuperscript{(20)} do \VUgras{vale} \textsuperscript{(25)} propriamente dito, em cujo
fundo se estende um bonito \VUgras{lago} \textsuperscript{(26)} que alimenta um
\VUgras{ribeiro} \textsuperscript{(24)} muito claro.

%%K t09-c1-11-1 p
11. --- Do outro lado do vale explora-se uma \VUgras{mina} \textsuperscript{(28)}
de ferro. Fui várias vezes ver descer os \VUgras{mineiros} pelo
\VUgras{poço}, e até entrei na \VUgras{galeria} onde passam o dia a
extrair o \VUgras{minério} que contém o \VUgras{metal}.

%%K t09-c1-12-1 p
12. --- Construiu-se uma importante \VUgras{fundição} junto à mina para
aproveitar já a seguir as riquezas que dela se tiram. O \VUgras{alto
forno} \textsuperscript{(29)}, aquecido com o \VUgras{carvão de pedra} que abunda
aqui, permite obter \VUgras{ferro fundido} de um modo rápido e barato.

%%K t09-c1-13-1 p
13. --- Não longe da mina abre-se uma \VUgras{pedreira} \textsuperscript{(30)}. Não
é nem \VUgras{granito}, nem \VUgras{pórfiro}, nem sequer \VUgras{grés} o que
o velho \VUgras{canteiro} arranca com a sua \VUgras{picareta}, mas simples
\VUgras{pedra de cantaria} para construir casas.

%%K t09-c1-14-1 p
14. --- Um dos mais belos adornos deste país é o \VUgras{abeto}
\textsuperscript{(31)}. Por toda a parte --- no fundo de um \VUgras{barranco}
\textsuperscript{(32)}, à beira de uma \VUgras{torrente} \textsuperscript{(33)}, junto a
um \VUgras{abismo} \textsuperscript{(34)} ou a um precipício ---, as suas finas
agulhas põem a sua nota verde.

%%K t09-tit-6 sub
\VUsaut{3.0mm}
\VUpk{10.2pt}{40}{Andar a pé.}
\VUsaut{3.0mm}

%%K t09-c1-15-1 p
15. --- Aproveito as minhas férias para dar longos passeios. Os
\VUgras{caminhos} \textsuperscript{(35)} que serpenteiam pela montanha permitem
percorrer, sem se dar conta da distância, vários \VUgras{quilómetros} e
muitas vezes várias \VUgras{léguas}.

%%K t09-c1-15-2 p
\VUgras{Marcos} \textsuperscript{(36)} e \VUgras{indicadores} \textsuperscript{(37)}
balizam os caminhos, nos quais se encontra muitas vezes um jovem
\VUgras{montanhês} \textsuperscript{(38)} com o \VUgras{alforge}
\textsuperscript{(40)} ao lado, que leva uma \VUgras{marmota} \textsuperscript{(39)}, ou
algum \VUgras{cigano} \textsuperscript{(41)} cujo \VUgras{barrete de pele}
\textsuperscript{(42)} contrasta de um modo singular com o seu velho
\VUgras{capuz} \textsuperscript{(43)} roto. Leva pela \VUgras{corrente} um
\VUgras{urso} \textsuperscript{(44)} amestrado.

%%K t09-tit-7 sub
\VUpk{10.2pt}{40}{Escalar montanhas.}
\VUsaut{3.0mm}

%%K t09-c1-16-1 p
16. --- Quase todos os dias vou com um \VUgras{guia} \textsuperscript{(48)} fazer
uma \VUgras{ascensão}, e estou muito orgulhoso quando, com a
\VUgras{mochila} \textsuperscript{(47)} às costas e o \VUgras{bastão alpino} na
mão, como um verdadeiro \VUgras{turista} \textsuperscript{(46)}, escalo as
\VUgras{rochas} \textsuperscript{(45)}.

%%K t09-c1-17-1 p
17. --- Do alto de uma montanha, os habitantes da planície não parecem
maiores do que formigas; um \VUgras{rebanho de ovelhas} \textsuperscript{(50)} que
pasta num \VUgras{pasto} \textsuperscript{(49)} parece um brinquedo de criança. Um
\VUgras{pinheiro} \textsuperscript{(51)}, ou mesmo uma \VUgras{cabana de madeira}
\textsuperscript{(52)}, mal são uma mancha na verdura.

%%K t09-c1-18-1 p
18. --- Durante estas excursões cruzamo-nos muitas vezes no
\VUgras{carreiro} \textsuperscript{(53)} com um \VUgras{caçador de camurças}
\textsuperscript{(54)} que leva ao ombro o animal que acaba de matar.

%%K t09-c1-19-1 p
19. --- Pus no meu herbário um ramo de \VUgras{feto} \textsuperscript{(55)} muito
curioso e um bonito raminho rosado de \VUgras{urze} \textsuperscript{(57)}, que
colhi à entrada de uma pequena \VUgras{gruta} \textsuperscript{(56)} no meu último
passeio.

%%K t09-tit-8 sub
\VUsaut{3.0mm}
\VUcentre{12.0pt}{0}{\textit{Segunda cena.}}
\VUsaut{3.0mm}

%%K t09-tit-9 sub
\VUpk{11.4pt}{40}{A Floresta. --- A Caça.}
\VUsaut{3.0mm}

%%K t09-c2-01-1 p
1. --- Ontem passei um dia delicioso na floresta. A atividade que reina
entre os animais e os \VUgras{insetos} \textsuperscript{(5)} que a povoam
interessou-me vivamente.

%%K t09-c2-01-2 p
A \VUgras{aranha} \textsuperscript{(1)} que tece a sua \VUgras{teia}
\textsuperscript{(2)} entre dois ramos para apanhar a \VUgras{mosca}
\textsuperscript{(3)} que passa, o \VUgras{caracol} \textsuperscript{(4)} que leva
penosamente a sua concha, o veado-volante que busca a sua presa, a
formiga diligente que se afadiga junto ao \VUgras{formigueiro}
\textsuperscript{(6)}: todo esse mundo pequeno ia, vinha e trabalhava com uma
animação extraordinária.

%%K t09-c2-02-1 p
2. --- Uma \VUgras{serpente} \textsuperscript{(7)}, que a princípio tomei por uma
\VUgras{víbora} mas que não era mais do que uma simples \VUgras{cobra},
estava acaçapada sob um tronco de árvore e de vez em quando deitava fora
a sua cabecinha fina, que sempre parecia disposta a morder. Diante de
mim, uma grossa \VUgras{lesma} \textsuperscript{(8)} rastejava pesadamente depois
de devorar um dos saborosos \VUgras{morangos silvestres} \textsuperscript{(11)}
que crescem aqui em abundância.

%%K t09-c2-03-1 p
3. --- O chão estava atapetado de \VUgras{flores da floresta}:
\VUgras{primaveras} \textsuperscript{(9)}, \VUgras{pervincas} \textsuperscript{(10)} e
muitas outras plantas que eu não conhecia floresciam entre as
\VUgras{moitas de erva} \textsuperscript{(12)}.

%%K t09-c2-04-1 p
4. --- Nesta região, a \VUgras{trufa} é quase tão comum como o
\VUgras{cogumelo} \textsuperscript{(14)}, e um \VUgras{porquinho}
\textsuperscript{(13)} de um guarda florestal procurava com ar guloso a ver se
encontrava alguma.

%%K t09-tit-10 sub
\VUsaut{3.0mm}
\VUpk{10.2pt}{40}{A caça furtiva.}
\VUsaut{3.0mm}

%%K t09-c2-05-1 p
5. --- Assisti ao \VUgras{fim} de uma cena que me divertiu muito. Graças
ao faro do seu \VUgras{cão basset} \textsuperscript{(15)}, o \VUgras{guarda
florestal} \textsuperscript{(16)} acabava de surpreender um \VUgras{caçador
furtivo} \textsuperscript{(22)} no momento em que este se dispunha a levar a
\VUgras{caça} \textsuperscript{(17)} que tinha morto. Preferindo abandonar a sua
presa a deixar-se prender, o furtivo tinha fugido, e \VUgras{lebre}
\textsuperscript{(18)}, \VUgras{corça} \textsuperscript{(19)}, \VUgras{corço}
\textsuperscript{(20)} e \VUgras{javali} \textsuperscript{(21)} ficavam sobre a erva,
para alegria do guarda.

%%K t09-c2-06-1 p
6. --- A variedade de árvores que a floresta encerra é assombrosa. As
\VUgras{faias} \textsuperscript{(23)} e as \VUgras{bétulas} \textsuperscript{(26)}, os
\VUgras{pinheiros}, que dão a \VUgras{resina} \textsuperscript{(27)}, os
\VUgras{carvalhos} \textsuperscript{(29)} e muitos mais entremeiam os seus ramos.

%%K t09-c2-06-2 p
A \VUgras{hera} \textsuperscript{(24)}, que se agarra ao tronco das árvores altas,
dá ao \VUgras{arvoredo alto} \textsuperscript{(25)} um verde brilhante que se une
agradavelmente ao tom mais claro e mais suave do \VUgras{musgo}
\textsuperscript{(31)} no qual o \VUgras{pica-pau-verde} \textsuperscript{(30)} busca
insetos.

%%K t09-tit-11 sub
\VUsaut{3.0mm}
\VUpk{10.2pt}{40}{Paisagem de floresta.}
\VUsaut{3.0mm}

%%K t09-c2-07-1 p
7. --- Os velhos carvalhos encantaram-me. As suas grossas \VUgras{raízes}
\textsuperscript{(32)} retorcidas, meio saídas da terra, dão-lhes uma esplêndida
aparência de força e de vigor, e pergunto-me como o \VUgras{proprietário}
\textsuperscript{(35)} se pode resignar a mandá-los abater e a entregá-los à
\VUgras{serra} \textsuperscript{(34)} do \VUgras{lenhador} \textsuperscript{(33)}. Os
pobres \VUgras{troncos} \textsuperscript{(36)}, despojados da sua \VUgras{casca}
\textsuperscript{(37)} ou fendidos pelo \VUgras{machado} \textsuperscript{(38)} do
\VUgras{jornaleiro} \textsuperscript{(39)}, dão-me pena.

%%K t09-c2-08-1 p
8. --- Perto dali, um velho \VUgras{carvoeiro} \textsuperscript{(41)} tinha
preparado com a sua \VUgras{pá} uma \VUgras{meda} \textsuperscript{(42)} de carvão,
e uma velha levava um \VUgras{feixe} \textsuperscript{(40)} de \VUgras{lenha
morta} feito com os raminhos das árvores abatidas.

%%K t09-tit-12 sub
\VUsaut{3.0mm}
\VUpk{10.2pt}{40}{A Caçada.}
\VUsaut{3.0mm}

%%K t09-c2-09-1 p
9. --- Queria ir descansar um pouco na \VUgras{casa do guarda florestal}
\textsuperscript{(46)}, mas quando cheguei junto ao \VUgras{tanque}
\textsuperscript{(43)}, no qual nada um belo \VUgras{cisne} \textsuperscript{(45)},
dei-me conta de que uma \VUgras{inundação} \textsuperscript{(44)} recente tinha
transformado o caminho num verdadeiro \VUgras{pântano}. Tive de dar uma
volta e, ao passar junto ao \VUgras{cedro} \textsuperscript{(49)}, aos
\VUgras{álamos} \textsuperscript{(48)} e ao \VUgras{olmo} \textsuperscript{(47)} que estão
na orla da floresta, vi sair de um \VUgras{souto} \textsuperscript{(54)} um
magnífico \VUgras{veado} \textsuperscript{(50)} perseguido por uma \VUgras{matilha}
\textsuperscript{(51)} encarniçada que a \VUgras{trompa} \textsuperscript{(53)} dos
\VUgras{monteiros a cavalo} \textsuperscript{(52)} ainda mais excitava. O pobre
animal estava esgotado. Veio cair moribundo a uns passos de mim.

%%K t09-c2-10-1 p
10. --- O filho do guarda florestal, atraído pelo ruído da
\VUgras{montaria}, saiu de casa e voltámos os dois à floresta. Tinha
visto uma \VUgras{pega} \textsuperscript{(56)} no ramo de um velho carvalho.
Pensava que o seu ninho devia estar escondido na \VUgras{folhagem}
\textsuperscript{(58)} e queria apanhá-lo.

%%K t09-c2-10-2 p
Ágil como um \VUgras{esquilo} \textsuperscript{(61)}, trepou de \VUgras{ramo}
\textsuperscript{(55)} em ramo, mas não encontrou nada e só conseguiu fazer voar
uma velha \VUgras{gralha} \textsuperscript{(60)} pousada num \VUgras{ramalho}
\textsuperscript{(57)}.
\VUsaut{4.0mm}
\VUfilet{22mm}
